A transição de modelos de Aprendizado Profundo avaliados em conjuntos de dados estáticos (\emph{offline benchmarks}) 
para sistemas de detecção de intrusão operando em redes de produção (\emph{online NIDS}) constitui um dos desafios mais 
complexos na engenharia de segurança contemporânea. Conforme discutido no Capítulo~\ref{cap:fundamentacao}, a literatura
acadêmica frequentemente restringe suas avaliações a arquivos de captura pré-gravados (\emph{Packet Capture} -- PCAP), 
processados em lote por ferramentas analíticas lentas como o \emph{CICFlowMeter}. Esse extrator, implementado 
primariamente em Java, não foi projetado para atuar como componente de baixa latência em linha de produção: sua 
arquitetura monolítica introduz contenção severa de memória volátil, atrasos imprevisíveis de coleta de lixo 
(\emph{garbage collection}) e ausência de suporte nativo a pipelines de tensores neurais acelerados por unidades de 
processamento gráfico (GPU) \cite{AOUINI2022108719}.

\section{Captura e Processamento de Tráfego}
\label{sec:captura_processamento}

Para que a arquitetura \textbf{ALF-MoE} possa exercer monitoramento e detecção perimetral em tempo hábil contra agentes 
adversariais, faz-se indispensável uma infraestrutura de telemetria capaz de inspecionar pacotes brutos diretamente da 
interface física de rede, agregá-los dinamicamente em fluxos bidirecionais segundo a quíntupla de fluxo
(\emph{5-tuple}), computar grandezas estatísticas e temporais em escala de microssegundos e despachá-las ao motor
neural sem degradar a vazão da rede. 

Neste capítulo, detalha-se a concepção, implementação e validação do pipeline de extração de tráfego baseado no 
\emph{framework} NFStream, o alinhamento semântico de suas características com o padrão canônico dos datasets do 
Canadian Institute for Cybersecurity (CIC) e o ambiente de execução em GPU para inferência contínua com auditoria 
forense.

\subsection{Extração com NFStream}
\label{subsec:extracao_nfstream}

O módulo de telemetria em tempo real do ecossistema ALF-MoE baseia-se no \emph{framework} NFStream
\cite{AOUINI2022108719}, selecionado pelo processamento em linguagem C de alta vazão, suporte nativo a
multiprocessamento e extensibilidade via plugins (\emph{User-Defined Plugins} -- UDPs). No pipeline, a classe
customizada \emph{FeatureExtractor} (derivada de \emph{NFPlugin}) atua como sensor passivo no interior dos processos de
trabalho (\emph{workers}), agregando grandezas estatísticas de rede diretamente dos pacotes brutos capturados na
interface física em modo promíscuo, sem acoplamento com tensores ou operações neurais.

Além dos contadores e momentos estatísticos agregados nativamente pela biblioteca, o extrator implementa duas rotinas
dedicadas de inspeção em nível de byte:
\begin{enumerate}
    \item \textbf{Codificação Compacta de Flags TCP (\emph{encode\_tcp\_flags}):} condensa os sinalizadores de controle
    do pacote em uma máscara binária de 8 bits (valores inteiros entre 0 e 255), viabilizando a identificação ágil de
    assinaturas anômalas e varreduras de portas (como \emph{TCP SYN}, \emph{NULL} e \emph{XMAS Scans});
    \item \textbf{Extração da Janela Inicial TCP (\emph{extract\_tcp\_window\_size}):} inspeciona o cabeçalho de
    transporte para recuperar o tamanho da janela inicial anunciada na abertura da conexão, diferenciando perfis
    legítimos de conexões sintetizadas por ferramentas de ataque. No caso de fluxos unidirecionais não respondidos, o
    valor sentinela de inicialização é automaticamente higienizado para zero no encerramento da conexão.
\end{enumerate}

O ciclo de vida dos fluxos no extrator é governado por dois temporizadores de expiração: o tempo limite de inatividade
($\tau_{\mathrm{idle}} = 120\,\mathrm{s}$), que consolida conexões após dois minutos de silêncio bilateral, e o tempo
limite ativo ($\tau_{\mathrm{active}} = 240\,\mathrm{s}$), que segmenta fluxos volumosos ou ataques contínuos de
negação de serviço (DoS/DDoS) a cada quatro minutos, impedindo a retenção indefinida de instâncias na memória de
telemetria.

Ao término da expiração, o vetor bruto de 76 atributos gerado pelo plugin é harmonizado com o pipeline analítico: os
intervalos temporais são escalonados de milissegundos para microssegundos ($\times 10^3$) para assegurar
compatibilidade com o padrão do CICFlowMeter, e os três atributos do sinalizador URG de variância nula são descartados,
resultando no vetor dimensional de 73 atributos ativos submetido à normalização robusta e logarítmica descritas nas
Seções~\ref{sec:pre-processamento} e \ref{subsec:normalizacao_codificacao}.

\subsection{Arquitetura de Concorrência e Contexto CUDA}
\label{subsec:concorrencia_cuda}

A separação de responsabilidades entre a extração de métricas e a classificação neural decorre do modelo de concorrência
do subsistema CUDA no ambiente Linux. Para viabilizar a ingestão contínua em redes de alta vazão, o NFStream adota uma
arquitetura orientada a multiprocessamento, delegando a captura em baixo nível a múltiplos processos de trabalho
(\emph{meters}) gerados pela chamada de sistema \emph{fork}.

Embora a memória virtual do processo pai seja duplicada para os processos filhos pelo mecanismo de cópia em gravação
(\emph{Copy-on-Write} ou CoW), os descritores de dispositivos aceleradores e o contexto dos drivers NVIDIA
(\emph{libcuda.so}) não são clonáveis nem compartilháveis entre processos bifurcados. Caso uma inferência neural fosse
invocada no retorno de encerramento do fluxo dentro de um processo filho, o driver interceptaria a chamada em um
processo sem contexto GPU válido, provocando falha na aplicação.

Para conciliar a aceleração por hardware na GPU dedicada com a estabilidade da telemetria de rede, o ecossistema
ALF-MoE implementa um padrão arquitetural Produtor-Consumidor desacoplado:
\begin{enumerate}
    \item \textbf{Produtor Assíncrono (Workers do NFStream):} os processos de medição do NFStream operam na captura de
    pacotes via \emph{libpcap} em modo promíscuo na interface física. No interior desses processos, o plugin
    \emph{FeatureExtractor} atua na CPU, agregando as 76 métricas nos objetos de fluxo e persistindo grandezas
    auxiliares na estrutura \emph{flow.udps}, sem interagir com tensores ou bibliotecas aceleradoras;
    
    \item \textbf{Consumidor Síncrono (Processo Principal com GPU):} o processo mestre itera sobre os fluxos
    expirados pelo temporizador do NFStream. Uma vez finalizado o ciclo do fluxo, o processo principal converte o
    objeto de medição em vetor de características, aplica a padronização estatística com o \emph{RobustScaler} na CPU
    e despacha a amostra diretamente para a memória da GPU, onde o grafo do ALF-MoE é executado sob um contexto CUDA
    exclusivo e persistente.
\end{enumerate}

Essa divisão garante estabilidade operacional contínua, prevenindo contenções de barramento e desacoplando o ritmo de
captura da rede do tempo de inferência neural do classificador.

\section{Setup de Inferência}
\label{sec:setup_inferencia}

Uma das premissas centrais desta pesquisa é comprovar a viabilidade operacional do modelo ALF-MoE em um arranjo 
experimental de inferência, mensurando a latência computacional de ponta a ponta e a robustez da auditoria forense em 
nível de sistema operacional.

\subsection{Arquitetura do Ambiente de Execução}
\label{subsec:arquitetura_execucao}

\begin{itemize}
    \item \emph{Equipamento:} Notebook Acer Nitro V15 (modelo ANV15-51);
    \item \emph{Processador:} Intel Core i7-13620H de 13ª geração (arquitetura híbrida Raptor Lake, dotado de 10
    núcleos físicos [6 \emph{Performance-cores} operando até 4,9 GHz e 4 \emph{Efficient-cores} até 3,6 GHz], 16
    threads lógicas e 24 MB de Intel Smart Cache);
    \item \emph{Memória Volátil:} 16 GB de memória RAM DDR5 operando em canal duplo com frequência de transferência
    de 5200 MT/s;
    \item \emph{Acelerador Gráfico Dedicado:} GPU NVIDIA GeForce RTX 4050 Laptop (arquitetura Ada Lovelace, 6 GB de
    memória de vídeo dedicada GDDR6, interface de 96 bits, 2560 núcleos CUDA e \emph{Tensor Cores} de 4ª geração);
    \item \emph{Sistema Operacional e Kernel:} Ubuntu 26.04.1 LTS de 64 bits, executando o kernel Linux versão
    7.0.0-38-generic;
    \item \emph{Pilha de Inteligência Artificial:} Python 3.11.9 (em ambiente virtual isolado \emph{.venv}) com
    bibliotecas TensorFlow 2.21.0 e Keras 3.15.1, operando com aceleração CUDA 12.5;
    \item \emph{Interface de Rede Local:} interface sem fio \emph{wlp63s0} conectada à sub-rede local 192.168.3.0/24
    com endereço IPv4 192.168.3.16.
\end{itemize}

A interface de rede sem fio da máquina (\emph{wlp63s0}) opera com a flag promíscua ativada pelo NFStream, assegurando
que todo o tráfego que transita pelo canal rádio direcionado ou originado pelo host seja capturado e espelhado para a
pilha de inspeção.

\subsection{Subsistema de Auditoria Forense e Registro SIEM (\emph{InferenceAuditor})}

O rastreamento e a persistência de eventos perimetrais são coordenados pelo módulo \emph{InferenceAuditor}
(implementado em \emph{src/inference.py}). Projetado para interoperabilidade com plataformas de gerenciamento de
informações e eventos de segurança (\emph{Security Information and Event Management} -- SIEM) e centros de operações de
segurança (SOC), o auditor opera com controle estrito de concorrência em ambiente multitarefa (\emph{threading.Lock}) e
gera dois fluxos paralelos de telemetria:
\begin{enumerate}
    \item \textbf{Log Operacional Legível Rotativo (\emph{.log}):} gerenciado pela classe \emph{RotatingFileHandler},
    com limite de $10\,\mathrm{MB}$ por arquivo e retenção de até cinco arquivos históricos de segurança
    (\emph{backup}), registrando mensagens operacionais formatadas para inspeção humana direta (com carimbo de
    data/hora, nível de severidade e PID do processo);
    \item \textbf{Trilha Forense Estruturada (\emph{.jsonl}):} registro atômico em formato JSON Lines, onde cada linha
    representa um documento JSON independente e autocontido com especificação temporal em formato ISO 8601 UTC
    (\emph{timestamp\_utc}).
\end{enumerate}

O evento principal de auditoria (\emph{FLOW\_PREDICTION}) grava a totalidade das variáveis que ditaram o julgamento
analítico do modelo para cada fluxo de rede inspecionado:
\begin{itemize}
    \item Metadados da quíntupla de fluxo (\emph{src\_ip}, \emph{src\_port}, \emph{dst\_ip}, \emph{dst\_port},
    \emph{protocol});
    \item Julgamento final do modelo ALF-MoE: classe predita (\emph{class}), probabilidade de confiança
    (\emph{confidence}) e sinalizador booleano de anomalia (\emph{is\_threat});
    \item Vetor pentadimensional de pesos da \emph{Gating Network}:
    \begin{equation}
    \boldsymbol{\alpha}_{\mathrm{GN}} = [\alpha_{\mathrm{DNN}}, \, \alpha_{\mathrm{CNN}}, \, \alpha_{\mathrm{GRU}},
    \, \alpha_{\mathrm{CAE}}, \, \alpha_{\mathrm{LSTM}}];
    \end{equation}
    \item Julgamentos e probabilidades individuais emitidos por cada um dos cinco especialistas neurais
    (\emph{experts\_predictions});
    \item Métrica de latência computacional estrita do \emph{forward pass} neural na GPU em milissegundos
    (\emph{latency\_ms});
    \item Ação de auditoria associada: indicação do registro analítico perimetral (\emph{LOG}) e estado operacional do
    evento.
\end{itemize}

Devido à natureza de concepção do dataset CSE-CIC-IDS2018 (Corrigido) -- que combina uma infraestrutura de rede
corporativa realista com múltiplos servidores de produção e estações de trabalho sob ataques diversificados, além do
saneamento metodológico das falhas de rotulação originais apontadas por \citeonline{liu2022error} --, este conjunto
foi selecionado como base exclusiva para a carga de pesos do modelo ALF-MoE integrado ao pipeline de inferência em
tempo real e auditoria forense do SIEM.

\subsection{Síntese Operacional do Pipeline de Inferência}
\label{subsec:sintese_pipeline_operacional}

A integração dos componentes detalhados neste capítulo consolida um ciclo operacional contínuo e determinístico para
inspeção de fluxos de rede e auditoria forense em tempo real:
\begin{enumerate}
    \item \textbf{Ingestão e Agregação de Tráfego:} os processos filhos do NFStream operam em modo promíscuo na
    interface física de rede, capturando pacotes brutos e agregando as 76 métricas canônicas do padrão CIC através do
    plugin customizado \emph{FeatureExtractor};
    \item \textbf{Desacoplamento e Expiração de Fluxos:} ao atingir o limiar de inatividade
    ($\tau_{\mathrm{idle}} = 120\,\mathrm{s}$) ou o tempo máximo ativo ($\tau_{\mathrm{active}} = 240\,\mathrm{s}$),
    o fluxo é encerrado nos workers em C/Cython e emitido de forma síncrona para o processo consumidor principal em
    Python, preservando a exclusividade do contexto CUDA no processo de inferência;
    \item \textbf{Conversão e Padronização:} o processo principal executa o alinhamento de unidades temporais
    ($\times 1000{,}0$ para compatibilidade com a escala em microssegundos do CICFlowMeter), expurga as três colunas
    de sinalizador URG de variância nula (reduzindo a entrada para 73 atributos ativos) e aplica o pré-processamento
    treinado (\emph{RobustScaler} e transformações logarítmicas $\ln(1+x)$);
    \item \textbf{Inferência Neural Acelerada:} o tensor resultante alimenta a arquitetura ALF-MoE na GPU, na qual a
    \emph{Gating Network} pondera os cinco especialistas neurais heterogêneos ($\mathcal{E}_{\mathrm{DNN}}$,
    $\mathcal{E}_{\mathrm{CNN}}$, $\mathcal{E}_{\mathrm{GRU}}$, $\mathcal{E}_{\mathrm{CAE}}$ e
    $\mathcal{E}_{\mathrm{LSTM}}$), consolidando a predição através da soma ponderada estrita e da cabeça classificadora
    profunda;
    \item \textbf{Auditoria Forense e Telemetria SIEM:} o \emph{InferenceAuditor} persiste o registro estruturado em
    JSONL com sincronização atômica (\emph{FLOW\_PREDICTION}), documentando a quíntupla do fluxo, a classe de ameaça
    identificada, a distribuição de confiança, os pesos atribuídos pela \emph{Gating Network}
    ($\boldsymbol{\alpha}_{\mathrm{GN}}$) e o tempo exato de latência computacional do modelo na GPU.
\end{enumerate}

O desempenho empírico desse arranjo experimental, abrangendo a acurácia de detecção nos benchmarks consolidados, a
latência de inferência por amostra e o comportamento dos especialistas frente aos diferentes vetores de ataque, é
apresentado e analisado exaustivamente no Capítulo~\ref{cap:resultados}.
